% answers_short.tex —— 2007 年·普通高等学校招生考试·天津卷·理科 · 简短答案速查
% 由 tools/make_short_answers.py 自 answers.tex 的 \gksolhead/\gkanswer 提取，
% 并与 2007/2007_数学真题_天津理_答案要点.md 逐题比对后生成，勿手改（改源 MD 后重跑）。

\begin{gknotice}
\textbf{使用说明}：本节只给最终答案便于快速核对，完整推导见解析版。答案为编写者独立推导，\textbf{非官方标准答案}。
\end{gknotice}

\gkpart{选择题}
\begin{center}\begin{tabular}{ccccccccccc}
\toprule
题号 & 1 & 2 & 3 & 4 & 5 & 6 & 7 & 8 & 9 & 10 \\
\midrule
答案 & C & B & A & D & C & D & B & B & A & A \\
\bottomrule
\end{tabular}\end{center}

\gkpart{填空题}
\begin{itemize}[leftmargin=2.6em]
  \item[\textbf{11.}] $2$。
  \item[\textbf{12.}] $14\pi$。
  \item[\textbf{13.}] $3$。
  \item[\textbf{14.}] $x+3y=0$。
  \item[\textbf{15.}] $-\frac23$。
  \item[\textbf{16.}] $390$。
\end{itemize}

\gkpart{解答题}
\begin{itemize}[leftmargin=2.6em]
  \item[\textbf{17.}] （1）最小正周期 $\pi$；（2）最小值 $-1$（$x=\frac{3\pi}{4}$），最大值 $\sqrt2$（$x=\frac{3\pi}{8}$）。
  \item[\textbf{18.}] （1）$\frac15$；（2）$\frac7{15}$；（3）$\xi$ 的分布列为 $P(\xi=0)=\frac15$，$P(\xi=1)=\frac7{15}$，$P(\xi=2)=\frac3{10}$，$P(\xi=3)=\frac1{30}$，$E\xi=\frac76$。
  \item[\textbf{19.}] （1）$CD\perp AE$（证明见解析）；（2）$PD\perp$ 平面 $ABE$（证明见解析）；（3）二面角 $A-PD-C$ 的大小为 $\arccos\frac{\sqrt2}{4}$（约 $69.3^\circ$）。
  \item[\textbf{20.}] （1）$6x+25y-32=0$；（2）$f'(x)=\dfrac{-2a\left(x-a\right)\left(x+\frac1a\right)}{(x^2+1)^2}$。当 $a>0$ 时，$f(x)$ 在 $\left(-\infty,-\frac1a\right)$、$(a,+\infty)$ 上递减，在 $\left(-\frac1a,a\right)$ 上递增，极小值 $f\left(-\frac1a\right)=-a^2$，极大值 $f(a)=1$。当 $a<0$ 时，$f(x)$ 在 $(-\infty,a)$、$\left(-\frac1a,+\infty\right)$ 上递增，在 $\left(a,-\frac1a\right)$ 上递减，极大值 $f(a)=1$，极小值 $f\left(-\frac1a\right)=-a^2$。
  \item[\textbf{21.}] （1）$a_n=(n-1)\lambda^n+2^n$；（2）当 $\lambda\ne1$ 时 $S_n=\dfrac{\lambda^2\left[1-n\lambda^{n-1}+(n-1)\lambda^n\right]}{(1-\lambda)^2}+2^{n+1}-2$；当 $\lambda=1$ 时 $S_n=\dfrac{n(n-1)}2+2^{n+1}-2$。
  \item[\textbf{22.}] （1）$a=\sqrt2\,b$（证明见解析）；（2）点 $D$ 的轨迹方程为 $x^2+y^2=\dfrac{2b^2}{3}$（即以 $O$ 为圆心、$\frac{\sqrt6}{3}b$ 为半径的圆）。
\end{itemize}
